\documentclass[11pt]{article}
\usepackage[T1]{fontenc}
\usepackage{lmodern,microtype}
\usepackage[a4paper,margin=24mm]{geometry}
\usepackage{amsmath,amssymb,amsthm}
\usepackage[hidelinks]{hyperref}
\newtheorem{theorem}{Theorem}
\newtheorem{lemma}[theorem]{Lemma}
\newcommand{\Z}{\mathbb Z}
\newcommand{\HH}{\mathcal H}
\setlength{\parindent}{0pt}
\setlength{\parskip}{4pt plus 1pt minus 1pt}
\emergencystretch=2em
\title{\bfseries The even supersingular branches of\\Swisher's (E.3) and (F.3)}
\author{}\date{}
\begin{document}
\maketitle\vspace{-3em}
\begin{abstract}
We prove a one-step congruence modulo $p^{3r}$ for two alternating
truncated hypergeometric sums. Its two-step consequence gives the
modulus $p^{3r-2}$ of Swisher's even branches at the original cutoffs,
including $p=3$ for (F.3). The proof uses a terminating Dougall sum,
quadratic interpolation and divisibility of every positive moment.
\end{abstract}
\section{The alternating sums}
Fix $d=3$ or $4$ and an odd prime $p\equiv-1\pmod d$. Let $a_r$ be
the element of $\{1,d-1\}$ congruent to $p^r$ modulo $d$.
For $a\in\{1,d-1\}$ and $N=a+dn$ with $n\ge0$, write
\[
 b_{a,k}=\frac{(a/d)_k}{k!},\qquad
 S_a(N)=\sum_{k=0}^n(-1)^k(2dk+a)b_{a,k}^3,\qquad
 T_r=S_{a_r}(p^r).
\]
Here $(x)_k$ is the rising factorial and $T_0=1$. The factor $2dk+a$
fixes the normalization for the companion family as well.
\begin{theorem}\label{thm}
For $r\ge1$ put $\epsilon_r=(-1)^{(p a_{r-1}-a_r)/d}$. Then
\[
 T_r\equiv p\epsilon_r T_{r-1}\pmod{p^{3r}}.
\]
Consequently, for every even $r\ge2$,
\[
 \boxed{S_1(p^r)\equiv p^2 S_1(p^{r-2})\pmod{p^{3r-2}}.}
\]
For $d=3$ this is the even supersingular (E.3) congruence, and for
$d=4$ it is the corresponding (F.3) congruence of Swisher's Table~2.
\end{theorem}
\section{A terminating evaluation}
Fix $a,N$ and abbreviate $b_k=b_{a,k}$, $H_i=a+di$, $E_i=di$. Put
\[
 S_n(y)=\sum_{k=0}^n(-1)^k(2dk+a)b_k^3
 \frac{\prod_{i=0}^{k-1}(1-y/H_i^2)}{\prod_{i=1}^k(1-y/E_i^2)}.
\]
At $H=H_h$, $0\le h\le n$, the terms beyond $h$ vanish and
\begin{equation}\label{node}
 S_n(H^2)=(-1)^hH.
\end{equation}
Indeed, put $\alpha=a/d$, $B=2\alpha+h$. The sum divided by $a$ is
the terminating \({}_4F_3(-1)\) with upper parameters
$\alpha,1+\alpha/2,B,-h$ and lower parameters
$\alpha/2,1+\alpha-B,1+\alpha+h$. Dougall's \({}_5F_4(1)\)
summation \cite{dlmf}, with one upper parameter $-h$ and another
tending to infinity, evaluates it as
$(1+\alpha)_h/(1+\alpha-B)_h=(-1)^h(\alpha+h)/\alpha$.
The limit is termwise in a finite sum.

\newpage
\section{Divisibility of all positive moments}
Define $\HH_a(N)=\{H_0,\ldots,H_n\}$,
$D_n(y)=\prod_{i=1}^n(1-y/E_i^2)$ and
$P_n(y)=\prod_{i=0}^n(1-y/H_i^2)$. Since $D_n S_n$ is a polynomial
of degree at most $n$, partial fractions and \eqref{node} give
\[
 F_n(y)=\frac{D_n(y)S_n(y)}{P_n(y)}
       =\sum_{H\in\HH_a(N)}\frac{\mu_{a,N}(H)}{1-y/H^2},
\]
where
\[
 \mu_{a,N}(H)=(-1)^{(H-a)/d} H
 \frac{\prod_{i=1}^n(1-H^2/E_i^2)}
      {\prod_{K\in\HH_a(N),\,K\ne H}(1-H^2/K^2)} .
\]
The nodes are distinct and none is a positive multiple of $d$.
Write $M_j(N)=\sum_H\mu_{a,N}(H)H^{2j}$, with $a$ fixed.
In particular $M_0(N)=F_n(0)=S_a(N)$.
\begin{lemma}\label{moments}
For each $j\ge0$ there is a polynomial $U_j\in\Z[1/2][X]$,
depending on $d,a$, such that
\[
 M_{j+1}(N)=(-1)^n(Nb_n)^3U_j(N),\qquad U_0=1.
\]
If $N=p^r$, all the moments $M_j(N)$ with $j\ge1$ belong to
$p^{3r}\Z_p$.
\end{lemma}
\begin{proof}
For $n\ge1$, separating the last term of $S_n$ gives
\[
 (1-y/N^2)F_n(y)
 =(1-y/(N-a)^2)F_{n-1}(y)+(-1)^n(2N-a)b_n^3.
\]
Expand $F_n(y)=-\sum_{j\ge1}M_j(N)y^{-j}$ at infinity.
The constant coefficient gives $M_1(N)=(-1)^n(Nb_n)^3$ by induction:
use $(N-d)b_{n-1}=(N-a)b_n$ and $M_1(a)=a^3$.
The coefficient of $y^{-j}$, $j\ge1$, gives
\[
 \frac{M_{j+1}(N)}{N^2}-M_j(N)=
 \frac{M_{j+1}(N-d)}{(N-a)^2}-M_j(N-d).
\]
After normalization this recurrence requires
\[
 XU_j(X)+(X-a)U_j(X-d)
 =X^3U_{j-1}(X)+(X-a)^3U_{j-1}(X-d).
\]
Over $R=\Z[1/2]$, take $U_{j-1}$ invariant under $X\mapsto a-d-X$.
The right side is odd about $a/2$. The operator on the left maps
polynomials even about $(a-d)/2$ to polynomials odd about $a/2$.
In the centered bases $(X-(a-d)/2)^{2i}$ and $(X-a/2)^{2i+1}$,
its matrix is triangular with diagonal $2$. It is invertible on
each matched finite degree filtration over $R$. Starting with $U_0=1$
therefore gives every $U_j$ in $R[X]$. At $X=a$ the recurrence gives
$U_j(a)=a^{2j}$, as required by the single node for $n=0$.
Induction in $j$ and then $n$ gives the asserted formula.
Since $-a/d\in\Z_p$ and $b_n=(-1)^n\binom{-a/d}{n}\in\Z_p$, it also
gives the divisibility. Binomial polynomials take integral values on
$\Z_p$ by continuity from the nonnegative integers.
\end{proof}

\newpage
\section{Transport through one prime-power level}
Put $N=p^r$, $M=p^{r-1}$, $a=a_r$, $b=a_{r-1}$ and
$e=(pb-a)/d$. Write $M=b+dm$; then $n=pm+e$ with $0\le e<p$.
The nodes divisible by $p$ are exactly $pJ$ for $J\in\HH_b(M)$.
The $p$-divisible denominator indices among $d,2d,\ldots,nd$ are
$pd,2pd,\ldots,mpd$. All these statements include $m=0$.

First remove the contribution of unit nodes. For $j_t=(p-1)p^t$,
$H^{2j_t}$ tends $p$-adically to $1$ if $p\nmid H$ and to $0$
otherwise. By Lemma~\ref{moments} and the finiteness of the node set,
\begin{equation}\label{units}
 \sum_{p\nmid H}\mu_{a,N}(H)=\lim_{t\to\infty}M_{j_t}(N)
         \in p^{3r}\Z_p.
\end{equation}
This limit is valid for rational weights of arbitrary individual valuation.

If $J=b+di$, then $(pJ-a)/d=pi+e$. Thus the ratio of the signed
node evaluations \eqref{node} is $p(-1)^e$. Cancelling the factors
with $p$-divisible indices in the weight formula yields
\[
 \mu_{a,N}(pJ)=p(-1)^e\mu_{b,M}(J)R(J),\qquad
 R(X)=\frac{\prod_{1\le i\le n,\,p\nmid E_i}(1-p^2X^2/E_i^2)}
 {\prod_{K\in\HH_a(N),\,p\nmid K}(1-p^2X^2/K^2)}.
\]
All indices in $R$ are units. Expanding its inverse factors geometrically,
\[
 R(X)=1+\sum_{j\ge1}c_jX^{2j},\qquad c_j\in p^{2j}\Z_p .
\]
The series converges on $\Z_p$. Summing over the finitely many $J$ and
using Lemma~\ref{moments} at the lower level shows that
\[
 p\sum_{J}\mu_{b,M}(J)(R(J)-1)
 =p\sum_{j\ge1}c_jM_j(M)\in
 p^{1+2+3(r-1)}\Z_p=p^{3r}\Z_p.
\]
Here $M_j(M)$ uses the residue class $b$. Together with \eqref{units}
this gives $T_r\equiv p(-1)^eT_{r-1}\pmod{p^{3r}}$.

At two consecutive levels the sign exponents add up to
$((pb-a)+(pa-b))/d=p-1$, since $a+b=d$. Their signs multiply to $1$.
The preceding step's error is multiplied by $p$, so the two-step
precision is at least $\min(3r,1+3(r-1))=3r-2$. At every even
level both endpoint sums have $a=1$, as in Theorem~\ref{thm}.
For $d=4,p=3,r=1$, $a=3$, $n=0$, and $T_1=3=pT_0$; all later
levels are covered by the same argument.

\small
\begin{thebibliography}{9}
\bibitem{swisher} H. Swisher, \emph{On the supercongruence conjectures of
van Hamme}, Research in the Mathematical Sciences 2 (2015), 18, Table 2.
\bibitem{gz} V. J. W. Guo and W. Zudilin, \emph{Dwork-type supercongruences
through a creative $q$-microscope}, J. Combin. Theory Ser. A 178 (2021),
105362, \S3.4.
\bibitem{guo} V. J. W. Guo, \emph{$q$-Analogues of two supercongruences in
the (E.3) and (F.3) conjectures of Swisher}, Rocky Mountain J. Math.
55 (2025), 147--154. The precise overlap with its theorems has not
been verified. This note makes no priority claim.
\bibitem{dlmf} NIST Digital Library of Mathematical Functions,
\href{https://dlmf.nist.gov/16.4.E9}{Eq.~16.4.9}, Rogers--Dougall
very-well-poised sum.
\end{thebibliography}
\end{document}
